% GENERATED FILE. Edit canonical lessons, then run npm run generate:books.
% canonical-source-hash: fnv1a64:06c8d22fc2204587
% canonical-lessons: ML-C305-niksepikkuka, ML-C305-pinvalikkuka, ML-C305-ceruka, ML-C305-kantumuttuka, ML-C305-sandarsikkuka

\chapter{To deposit, To withdraw, To join, To meet, To visit}
\label{ch:ml-to-deposit-to-withdraw-to-join-to-meet-to-visit}
\hlchaptermodality{305}

\begin{chapteropening}
\textbf{By the end of this chapter:} \emph{I can say to deposit, to withdraw, to join, to meet and to visit.}

Complete the last lesson of chapter 305: I can say to deposit, to withdraw, to join, to meet and to visit.
\end{chapteropening}

\section[\texorpdfstring{\ml{നിക്ഷേപിക്കുക} (nikṣēpikkuka)}{നിക്ഷേപിക്കുക (nikṣēpikkuka)}]{\ml{നിക്ഷേപിക്കുക} (nikṣēpikkuka) — to deposit}

\label{lesson:ML-C305-niksepikkuka}



\begin{quote}
Before the new one: say the Malayalam for to swell, then the Malayalam for to itch.
\end{quote}

\subsection*{You'll want to know: \ml{നിക്ഷേപിക്കുക}}
\textbf{\ml{നിക്ഷേപിക്കുക}} — \emph{nikṣēpikkuka} — \textquotedblleft{}to deposit\textquotedblright{}.

\begin{grammarlens}[title={What you've built}]
One more word for this chapter.
\end{grammarlens}

\subsection*{Your turn}
\begin{itemize}
\raggedright
  \item \emph{Say it:} \emph{nikṣēpikkuka}
  \item \emph{Say it:} \emph{nikṣēpikkuka}, once more
  \item \emph{Say it:} say \emph{nikṣēpikkuka}
  \item \emph{Recall:} say the Malayalam for to swell, then the Malayalam for to itch, then say \emph{nikṣēpikkuka} again
\end{itemize}

\begin{tcolorbox}[breakable,colback=teal!4,colframe=teal!35!black,title={Before you move on}]
What does \ml{നിക്ഷേപിക്കുക} mean? (\textquotedblleft{}To deposit\textquotedblright{}.) Say it once more.
\end{tcolorbox}
\section[\texorpdfstring{\ml{പിൻവലിക്കുക} (pinvalikkuka)}{പിൻവലിക്കുക (pinvalikkuka)}]{\ml{പിൻവലിക്കുക} (pinvalikkuka) — to withdraw}

\label{lesson:ML-C305-pinvalikkuka}



\begin{quote}
Before the new one: say the Malayalam for to itch, then the Malayalam for to deposit.
\end{quote}

\subsection*{You'll want to know: \ml{പിൻവലിക്കുക}}
\textbf{\ml{പിൻവലിക്കുക}} — \emph{pinvalikkuka} — \textquotedblleft{}to withdraw\textquotedblright{}.

\begin{grammarlens}[title={What you've built}]
One more word for this chapter.
\end{grammarlens}

\subsection*{Your turn}
\begin{itemize}
\raggedright
  \item \emph{Say it:} \emph{pinvalikkuka}
  \item \emph{Say it:} \emph{pinvalikkuka}, once more
  \item \emph{Say it:} say \emph{pinvalikkuka}
  \item \emph{Recall:} say the Malayalam for to itch, then the Malayalam for to deposit, then say \emph{pinvalikkuka} again
\end{itemize}

\begin{tcolorbox}[breakable,colback=teal!4,colframe=teal!35!black,title={Before you move on}]
What does \ml{പിൻവലിക്കുക} mean? (\textquotedblleft{}To withdraw\textquotedblright{}.) Say it once more.
\end{tcolorbox}
\section[\texorpdfstring{\ml{ചേരുക} (cēruka)}{ചേരുക (cēruka)}]{\ml{ചേരുക} (cēruka) — to join}

\label{lesson:ML-C305-ceruka}



\begin{quote}
Before the new one: say the Malayalam for to deposit, then the Malayalam for to withdraw.
\end{quote}

\subsection*{You'll want to know: \ml{ചേരുക}}
\textbf{\ml{ചേരുക}} — \emph{cēruka} — \textquotedblleft{}to join\textquotedblright{}.

\begin{grammarlens}[title={What you've built}]
One more word for this chapter.
\end{grammarlens}

\subsection*{Your turn}
\begin{itemize}
\raggedright
  \item \emph{Say it:} \emph{cēruka}
  \item \emph{Say it:} \emph{cēruka}, once more
  \item \emph{Say it:} say \emph{cēruka}
  \item \emph{Recall:} say the Malayalam for to deposit, then the Malayalam for to withdraw, then say \emph{cēruka} again
\end{itemize}

\begin{tcolorbox}[breakable,colback=teal!4,colframe=teal!35!black,title={Before you move on}]
What does \ml{ചേരുക} mean? (\textquotedblleft{}To join\textquotedblright{}.) Say it once more.
\end{tcolorbox}
\section[\texorpdfstring{\ml{കണ്ടുമുട്ടുക} (kaṇṭumuṭṭuka)}{കണ്ടുമുട്ടുക (kaṇṭumuṭṭuka)}]{\ml{കണ്ടുമുട്ടുക} (kaṇṭumuṭṭuka) — to meet}

\label{lesson:ML-C305-kantumuttuka}



\begin{quote}
Before the new one: say the Malayalam for to withdraw, then the Malayalam for to join.
\end{quote}

\subsection*{You'll want to know: \ml{കണ്ടുമുട്ടുക}}
\textbf{\ml{കണ്ടുമുട്ടുക}} — \emph{kaṇṭumuṭṭuka} — \textquotedblleft{}to meet\textquotedblright{}.

\begin{grammarlens}[title={What you've built}]
One more word for this chapter.
\end{grammarlens}

\subsection*{Your turn}
\begin{itemize}
\raggedright
  \item \emph{Say it:} \emph{kaṇṭumuṭṭuka}
  \item \emph{Say it:} \emph{kaṇṭumuṭṭuka}, once more
  \item \emph{Say it:} say \emph{kaṇṭumuṭṭuka}
  \item \emph{Recall:} say the Malayalam for to withdraw, then the Malayalam for to join, then say \emph{kaṇṭumuṭṭuka} again
\end{itemize}

\begin{tcolorbox}[breakable,colback=teal!4,colframe=teal!35!black,title={Before you move on}]
What does \ml{കണ്ടുമുട്ടുക} mean? (\textquotedblleft{}To meet\textquotedblright{}.) Say it once more.
\end{tcolorbox}
\section[\texorpdfstring{\ml{സന്ദർശിക്കുക} (sandarśikkuka)}{സന്ദർശിക്കുക (sandarśikkuka)}]{\ml{സന്ദർശിക്കുക} (sandarśikkuka) — to visit}

\label{lesson:ML-C305-sandarsikkuka}



\begin{quote}
Before the new one: say the Malayalam for to join, then the Malayalam for to meet.
\end{quote}

\subsection*{You'll want to know: \ml{സന്ദർശിക്കുക}}
\textbf{\ml{സന്ദർശിക്കുക}} — \emph{sandarśikkuka} — \textquotedblleft{}to visit\textquotedblright{}.

\begin{grammarlens}[title={What you've built}]
One more word for this chapter.
\end{grammarlens}

\subsection*{Your turn}
\begin{itemize}
\raggedright
  \item \emph{Say it:} \emph{sandarśikkuka}
  \item \emph{Say it:} \emph{sandarśikkuka}, once more
  \item \emph{Say it:} say \emph{sandarśikkuka}
  \item \emph{Recall:} say the Malayalam for to join, then the Malayalam for to meet, then say \emph{sandarśikkuka} again
\end{itemize}

\begin{tcolorbox}[breakable,colback=teal!4,colframe=teal!35!black,title={Before you move on}]
What does \ml{സന്ദർശിക്കുക} mean? (\textquotedblleft{}To visit\textquotedblright{}.) Say it once more.
\end{tcolorbox}
